\section{The actual original germ ring and its full formal completion}

Let $R$ be the ring of original holomorphic germs at the origin of
$\mathbb C^n$, let $\mathfrak m$ be its actual maximal ideal, and let
$A=\mathbb C[x_1,\ldots,x_n]$ with variable ideal $\mathfrak v$.
Polynomial evaluation is at the original coordinate germs. The constructions
below use the already established genuine finite Taylor representatives;
no convergent realization of an arbitrary formal series is asserted.

\begin{theorem}[Actual polynomial jets are the whole original local jet rings]
\label{hir:actual-original-polynomial-jet-ring-equivalences}
\uses{hir:actual-original-polynomial-low-coefficient-separation,hir:actual-original-canonical-finite-taylor-polynomial,hir:actual-original-formal-variable-power-extension}
\lean{Hironaka.HironakaActualOriginalAdicCompletion.originalPolynomial_mem_maximal_pow_iff,Hironaka.HironakaActualOriginalAdicCompletion.originalPolynomialJetEquiv,Hironaka.HironakaActualOriginalAdicCompletion.originalPolynomialJetEquiv_transition}
\leanok
For every $q\geq0$, genuine polynomial evaluation induces a ring
isomorphism
\[
 e_q:A/\mathfrak v^q\;\xrightarrow{\;\sim\;}R/\mathfrak m^q,
 \qquad [P]\longmapsto[P(x)].
\]
These isomorphisms commute with every quotient reduction $q\to p$ for
$p\leq q$. The whole preimage of $\mathfrak m^q$ is exactly
$\mathfrak v^q$.
\end{theorem}
\begin{proof}
The actual Taylor homomorphism takes $P(x)$ to the same polynomial as a
formal series. Membership of an original germ in $\mathfrak m^q$ is
equivalent to the vanishing of every Taylor coefficient of total degree
less than $q$. Applied to $P(x)$, these are precisely the coefficients of
$P$; their vanishing is exactly membership in $\mathfrak v^q$. This proves
both whole ideal directions and makes the indicated quotient homomorphism
well defined and injective.

For $q>0$, the genuine finite Taylor polynomial of order $q-1$ of any
original germ $g$ has remainder in $\mathfrak m^q$. Its polynomial class
therefore maps to the entire class of $g$, so the map is surjective.
For $q=0$ both quotients are the zero ring and the same conclusion holds.
Every quotient element has a polynomial representative, and both
composites around a reduction square send that representative to the
class of the same actual evaluated polynomial modulo $\mathfrak m^p$.
Thus the reduction squares commute, without choosing compatible
representatives as additional data.
\end{proof}

\begin{theorem}[The genuine maximal-adic completion is the actual formal Taylor ring]
\label{hir:actual-original-maximal-adic-formal-ring-equivalence}
\uses{hir:actual-original-polynomial-jet-ring-equivalences,hir:actual-original-all-taylor-coefficients-faithful}
\lean{Hironaka.HironakaActualOriginalAdicCompletion.actualOriginalCompletionEquiv,Hironaka.HironakaActualOriginalAdicCompletion.actualOriginalCompletionEquiv_of}
\leanok
There is an internally constructed ring isomorphism
\[
 \widehat R_{\mathfrak m}\;\xrightarrow{\;\sim\;}
 \mathbb C[[x_1,\ldots,x_n]]
\]
whose composite with the original completion map $R\to\widehat R$ is
the genuine original Taylor homomorphism on every original germ.
\end{theorem}
\begin{proof}
An element of either adic completion is a family of classes, one in each
power quotient, coherent under every quotient reduction. Apply $e_q$ to
each component of a coherent polynomial-jet family. The reduction squares
of the first theorem show that the resulting original-jet family is
coherent. Applying $e_q^{-1}$ componentwise gives the reverse construction;
the same reduction squares, after composing with the two inverse maps,
prove its coherence. Both maps preserve zero, one, addition and
multiplication componentwise. Their two composites are identities in
every quotient, hence identities of the coherent families. This identifies
the original maximal-adic completion with the variable-adic completion of
the actual polynomial ring.

Use the canonical formal-series isomorphism with the latter polynomial
completion, given in degree $q$ by truncation to monomials of total degree
less than $q$. To check the required original embedding identity, fix an
original germ $g$ and $q>0$. Its finite Taylor polynomial of order $q-1$
represents its entire original class. Its coefficients of degree less than
$q$ are exactly those of the genuine original Taylor series, so its
polynomial quotient class is the truncation class of that series. For
$q=0$ both classes are zero. Thus the two coherent polynomial families
agree at every degree, and the formal-series isomorphism gives exactly the
original Taylor series of $g$. This argument establishes an isomorphism of
completions, not convergence of arbitrary formal series in the original
analytic germ ring.
\end{proof}
